\section{Experimental setup}
\label{sec:setup}

\subsection{Research questions}
\label{sec:rqs}

The evaluation is organized around six research questions:
\begin{itemize}
\item \textbf{RQ1 (reliability)}: Which deployment methods actually
keep their contract --- verified directly on the \emph{exact risk over
a fixed finite population} across repeated calibration draws and
across contract levels --- and does the guarantee extend to
identifiable query slices?
\item \textbf{RQ2 (cost-efficiency)}: At what cost does the certified
policy meet the contract, relative to fixed pipelines, uncertified
tuners, and hindsight oracles --- including joint multi-dimensional
contracts and strict low-$\alpha$ regimes with certified deferral?
\item \textbf{RQ3 (risk vouchers)}: Do vouchers achieve their nominal
coverage over repeated calibration draws, how loose is assumption-free
composition versus direct certification, and does selectivity
conditioning expose real heterogeneity?
\item \textbf{RQ4 (knowledge drift)}: Are workload and index-staleness
drifts detected, how fast, under full and subsampled auditing, with
what lifetime guarantee across repeated alarms --- and what does the
auditing itself cost?
\item \textbf{RQ5 (generalization)}: Do the guarantees and the savings
survive a change of LLM family, model developer, serving backend, and
cost metric --- including a locally served open-weight family metered
in GPU-seconds?
\item \textbf{RQ6 (planning and scalability)}: Does cardinality-driven
knowledge-source selection beat fixed acquisition orders, and how does
certification scale to $10^5$-candidate plan spaces?
\end{itemize}

\subsection{Datasets and contracts}
\label{sec:datasets}

We evaluate on four heterogeneous benchmarks with fully metered
execution.
(A) \textbf{HybridQA} \citep{chen2020hybridqa}: multi-hop QA over
Wikipedia tables joined with entity passages; sources are a structured
row-selection operator and sparse/dense passage indexes over 440K
passages; quality loss $\mathbf{1}[\text{token-F1} < 0.5]$; splits
train/cal/test $= 5000/3000/3466$ (test is the entire official
development release).
(B) \textbf{CRAG} \citep{yang2024cragbench}: web QA with the official
mock knowledge-graph API (five domains, dynamism labels
static/slow/fast/real-time); evidence chunks carry crawl ages; answers
are graded by an LLM judge with the gold answer following CRAG's
protocol; loss is correctness ($\mathbf{1}[\text{not correct}]$) or
hallucination ($\mathbf{1}[\text{answered and wrong}]$); splits
$700/1000/1006$ covering the full public development pool.
(C) \textbf{ASQA} \citep{stelmakh2022asqa} under the official ALCE
evaluation \citep{gao2023alce}: long-form answers with citations over a
pooled GTR corpus (89K passages), citation recall scored by the
official TRUE-NLI entailment model \citep{honovich2022true}; loss
$\mathbf{1}[\text{citation recall} < 50]$; splits $150/350/448$
covering the full release.
(D) \textbf{QAMPARI} \citep{amouyal2022qampari} under the same ALCE
protocol: list answers over an 86K-passage corpus; loss
$\mathbf{1}[\text{citation recall} < 50]$; splits $150/350/500$.
QAMPARI is the adversarial regime for citation contracts --- even the
strongest plan violates on 61\% of queries --- so it probes
certification exactly where point estimates are noisiest.

These scales reflect the operating regime of a deployed knowledge
service that recalibrates periodically on a few hundred to a few
thousand audited queries: at
$n_{\mathrm{cal}} \in \{350, 1000, 3000\}$ and $\delta = 0.1$, the
Hoeffding--Bentkus certificate half-width at mid-range $\alpha$ is
roughly $0.075/0.04/0.025$ --- tight enough that certificates
discriminate among ladder rungs rather than collapsing to the strongest
plan. For ASQA, the smallest track, no single-split number is
load-bearing: the RQ1 protocol resamples 1000 calibration/test
partitions from the pooled 798 queries.

\subsection{Pipelines, ladders, and cost metering}
\label{sec:ladders}

Each benchmark has a four-rung cost-ascending ladder generated by the
rewrites of Section~\ref{sec:rewrites}: from sparse-only retrieval with
a flash-tier model up to hybrid reciprocal-rank-fusion retrieval,
cross-encoder reranking, NLI verification, knowledge-graph-API evidence
(CRAG), citation repair (ASQA, QAMPARI), and a max-tier model. The
physical layer uses DuckDB for structured sources, BM25 for sparse
retrieval, BGE embeddings with FAISS for dense retrieval, a
bge-reranker-v2-m3 cross-encoder, and DeBERTa-based NLI verification;
local models run on two RTX~4090 GPUs. All LLM executions are real API
calls with exact token metering; reported cost is LLM spend in
milli-CNY per query. A query stopping at rung $j$ is charged rungs
$0..j$ cumulatively for escalation policies and only rung $j$ for jump
policies. All LLM and API calls flow through a thread-safe cache keyed
by (model, messages, decoding parameters), making every experiment
reproducible and token-exact. RQ5 additionally swaps the generator
column of the HybridQA ladder across four model families from four
\emph{model developers} --- Alibaba's Qwen (flash/plus/max),
DeepSeek (v4-flash/v3.2/v4-pro), and Zhipu's GLM (4.7/5/5.2), plus
Google's open-weight Gemma~3 (1b/4b/12b/27b) --- served through two
\emph{serving backends}: the three commercial families are hosted on
one API platform (DashScope) and token-metered at each developer's
list prices, while Gemma~3 is served locally by Ollama and metered in
measured GPU-seconds priced at an RTX-4090 rental rate. We keep the
two notions distinct throughout: the experiment varies the model
family (and with it the price structure), not the number of API
platforms. All families run over shared materialized retrieval
(Section~\ref{sec:planner}), so hosted and local plans share one cost
currency.

\subsection{Baselines}
\label{sec:baselines}

We compare against representative members of every deployment paradigm,
all consuming identical execution matrices so comparisons are exact.
\emph{Fixed plans}: each rung as a static pipeline. \emph{Complexity
router}: Adaptive-RAG \citep{jeong2024adaptiverag} adapted to our
ladders --- a classifier trained to predict the cheapest correct rung.
\emph{Utility router}: per-query
$\arg\max_j \hat{s}_j - \lambda \bar c_j$ with $\lambda$ tuned to the
constraint on tuning data
\citep{ding2024hybridllm,hu2024routerbench}. \emph{BO-tuned}:
Optuna \citep{akiba2019optuna} tree-structured Parzen estimation over
escalation thresholds, minimizing cost subject to empirical risk
$\le \alpha$ on tuning data. \emph{Abacus} \citep{russo2025abacus}:
workload-level constrained selection with UCB sampling, choosing the
cheapest rung whose estimated risk clears $\alpha$.
\emph{Empirical-no-cert}: our own ladder with the cheapest thresholds
whose empirical calibration risk is $\le\alpha$ --- i.e., \system\
minus the statistical test (ablation). We report two hindsight
skylines with different targets: the \emph{zero-loss oracle} picks the
per-query cheapest rung that incurs no loss (a stricter objective than
any contract, hence not cost-comparable to contract-feasible methods),
and the \emph{matched oracle} is the cheapest per-query rung
assignment whose \emph{test-set} mean loss is $\le \alpha$ --- the
hindsight optimum of the same constrained problem \system\ solves, and
the fair lower bound for its cost.

\subsection{Evaluation protocol}
\label{sec:protocol}

For reliability (RQ1) we verify the theory's actual claim --- a bound
on the probability that the deployed policy's \emph{true} risk exceeds
$\alpha$ --- directly, without the binomial noise of a small held-out
sample. We treat the full pooled execution matrix as a fixed finite
population, repeat 1000 times: draw a random calibration subset, let
every method select and tune on it (scorers and routers use the fixed
train split), and evaluate the selected policy's \emph{exact} risk on
the entire population. We report
$\Pr(\text{population risk} > \alpha)$ --- the exact finite-population
analogue of the $\delta$-guarantee, with $\delta = 0.1$ throughout ---
alongside the conventional held-out estimate; at 1000 draws the
two-sided 95\% interval on a reported rate of $0.05$ is $\pm 0.014$.
Because calibration sets are drawn \emph{without replacement} from the
pooled matrix while the certificate's p-value is derived for i.i.d.\
draws, Appendix~\ref{app:finitepop} proves the formal bridge: the
hypergeometric
test is exact for this protocol, the shipped p-value is dominated by it
at every configuration used here (so it stays valid, merely
conservative), and re-certifying every draw with the exact
finite-population p-value leaves all conclusions unchanged while
certifying cheaper policies.
Throughout the paper we keep three quantities distinct: (i)
\emph{certification} is an event on the calibration draw; (ii) the
guarantee it purchases is on the \emph{population} risk, with failure
probability $\le \delta$ over draws; (iii) the \emph{test-split
empirical risk} of a single draw is a noisy estimate of (ii) and may
exceed $\alpha$ without any guarantee being violated. Tables that show
a single draw are labelled accordingly. Cost-efficiency (RQ2) fixes
representative contract levels per benchmark and reports realized test
risk and mean cost. The remaining protocols are described with their
results in Section~\ref{sec:results}.
